\begin{table}[H]
\centering\small
\caption{Seven radiomics feature families extracted from the 3D MRI volume.}
\label{tab:radiomics-groups}
\begin{tabularx}{\linewidth}{@{}L{2.7cm}r Y Y@{}}
\toprule
Family & Count & Descriptive information & Representative examples \\
\midrule
Shape & 14 & Geometric properties of the whole-brain mask & MeshVolume, SurfaceArea \\
First-order & 18 & Voxel intensity distribution within the mask & Mean, Variance, Entropy \\
GLCM & 24 & Gray-level co-occurrence relationships across voxel pairs & Contrast, Correlation \\
GLRLM & 16 & Run lengths of consecutive voxels sharing a gray level & ShortRunEmphasis, LongRunEmphasis \\
GLSZM & 16 & Sizes of connected homogeneous gray-level zones & SmallAreaEmphasis, LargeAreaEmphasis \\
GLDM & 14 & Gray-level dependency between a voxel and its neighborhood & DependenceEntropy, DependenceVariance \\
NGTDM & 5 & Intensity difference between a voxel and its neighbors & Coarseness, Contrast, Busyness, Complexity, Strength \\
\midrule
Total & 107 & 104 retained after removing Maximum, Minimum, and Range & PCA retains 16 components \\
\bottomrule
\end{tabularx}
\caption*{\textit{Note.} Definitions and representative features follow PyRadiomics~\cite{pyradiomicsfeatures}. Feature counts reflect the extraction catalog of this thesis.}
\end{table}
